\section{Discussion}
SpecGate's benefit comes almost entirely from deciding \emph{when not to prefetch}, not from a richer prediction model. This has two practical implications. First, the mechanism is cheap: a 2 KB PTT with 3-bit counters is a rounding error next to a modern L2, and the gate itself is a single comparator on the lookup path, adding no measurable cycles to the miss handler in our RTL-level latency model. Second, the same trigger table could be paired with a more sophisticated predictor (e.g., a multi-delta or context predictor) without changing the gating logic, since confidence tracking is orthogonal to how the successor address is computed.

The threshold $\theta$ is a coverage/accuracy trade-off: lowering it increases coverage at the cost of more wasted prefetches, while raising it suppresses noise but delays how quickly a genuinely stable trigger starts issuing. We swept $\theta \in \{0.3, 0.4, 0.5, 0.6, 0.7\}$ and found $\theta = 0.5$ within 1.2\% of the best geomean speedup for every kernel we tested, which is why we report it as the default; workloads with faster-changing trigger populations may prefer a lower threshold to reduce the ramp-up penalty on new entries. A second limitation is capacity: at 512 entries, the PTT itself thrashes on hash-join's largest table sizes, and we expect diminishing returns from further increasing PTT size without also improving its replacement policy, which we leave to future work. Finally, our evaluation is single-core; extending confidence gating to a shared, multi-core PTT would require partitioning or per-core confidence tracking to avoid one core's noisy triggers suppressing another core's reliable ones.
